Let \((\mathsf Y,\mathcal Y)\), \(e\), and \(\mathbf Y\) be furnished by Assumption~\texttt{ass:potential-outcome-path-typing}. Because \(\mathsf Y\) and \([0,1]^d\) are standard Borel spaces, there is a regular conditional probability kernel
\[
K_x=\mathcal L(\mathbf Y\mid X=x).
\]
Define
\[
\overline m_P(a,x)=\int_{\mathsf Y}e(y,a)K_x(dy).
\]
The evaluation map is bounded and jointly measurable, and \(x\mapsto K_x\) is a probability kernel. Kernel integration, first for simple functions and then by bounded measurable approximation, shows that \((a,x)\mapsto\overline m_P(a,x)\) is jointly Borel. For every fixed \(a\), it is a version of \(\mathbb E[Y(a)\mid X=x]\).

Process-level conditional exchangeability states that a conditional law of \(\mathbf Y\) given \((A,X)\) is \(K_X\). Applying this kernel identity to the bounded jointly measurable function \((y,a,x)\mapsto e(y,a)\) gives
\[
\mathbb E[e(\mathbf Y,A)\mid A,X]
=\int_{\mathsf Y}e(y,A)K_X(dy)
=\overline m_P(A,X)
\quad\text{almost surely}.
\]
The corrected path-typing identity holds at every sample point and for every dose. Substituting \(a=A\) therefore gives \(Y(A)=e(\mathbf Y,A)\) immediately, with no union of dose-specific null sets. Consistency gives \(Y=Y(A)\) almost surely. Hence
\[
\mathbb E[Y\mid A,X]=\overline m_P(A,X)
\quad\text{almost surely}.
\]
The designated \(\mu_P(A,X)\) is another version of this conditional mean, so
\[
\mu_P(A,X)=\overline m_P(A,X)
\quad\text{almost surely}.
\]
The joint law of \((A,X)\) has density \(\pi_P(a\mid x)\) relative to \(da\,dP_X(x)\). Consequently,
\[
0=\int_{[0,1]^d}\int_0^1
\lvert\mu_P(a,x)-\overline m_P(a,x)\rvert
\pi_P(a\mid x)\,da\,dP_X(x).
\]
Uniform positivity \(\pi_P\ge c>0\) upgrades this equality to
\[
\mu_P(a,x)=\overline m_P(a,x)
\quad\text{for }da\,dP_X(x)\text{-almost every }(a,x). \tag{1}
\]

Now let \(m_P\) be the jointly Borel counterfactual conditional-mean version in the proposition. For every fixed \(a\), both \(m_P(a,\cdot)\) and \(\overline m_P(a,\cdot)\) are versions of \(\mathbb E[Y(a)\mid X=\cdot]\), and hence they agree \(P_X\)-almost everywhere. Tonelli's theorem and joint measurability then imply
\[
m_P(a,x)=\overline m_P(a,x)
\quad\text{for }da\,dP_X(x)\text{-almost every }(a,x). \tag{2}
\]
Combining (1) and (2) proves the derived bridge
\[
\mu_P(a,x)=m_P(a,x)
\quad\text{for }da\,dP_X(x)\text{-almost every }(a,x). \tag{3}
\]
By Fubini's theorem, for \(P_X\)-almost every \(x\), equality in (3) holds for Lebesgue-almost every \(a\), the section \(a\mapsto m_P(a,x)\) is continuous at \(t_0\), and the treatment section \(a\mapsto\mu_P(a,x)\) is Hölder and therefore continuous. Fix such an \(x\). Every neighbourhood of the interior point \(t_0\) has positive Lebesgue measure, so one may choose equality points \(a_k\to t_0\). Passing to the limit on both sides gives
\[
\mu_P(t_0,x)=m_P(t_0,x).
\]
Thus this equality holds for \(P_X\)-almost every \(x\). Since \(m_P(t_0,X)\) is a version of \(\mathbb E[Y(t_0)\mid X]\), the density identity \(dP_X(x)=p_X(x)\,dx\) and the tower property yield
\[
\theta_P(t_0)
=\int\mu_P(t_0,x)p_X(x)\,dx
=\int m_P(t_0,x)\,dP_X(x)
=\mathbb E[Y(t_0)].
\]
Proposition~\texttt{prop:estimand-version-invariant} ensures that the observed-data expression does not depend on the admissible designated versions.

It remains to establish the singleton-dose obstruction inside the stated path typing. Fix the observed law
\[
X\sim\operatorname{Unif}([0,1]^d),\qquad
A\mid X\sim\operatorname{Unif}([0,1]),\qquad Y=0,
\]
with \(p_X(x)=1\), \(\pi_P(a\mid x)=1\), and \(\mu_P(a,x)=0\). The same checks as in Proposition~\texttt{prop:class-nonempty} put this law in every admissible observed-data class. Choose \(0<\eta\le M_Y\). Let \(\mathsf Y=\{y_0,y_1\}\) with the discrete sigma-field, and define
\[
e(y_0,a)=0,
\qquad
e(y_1,a)=\eta\mathbf 1\{a=t_0\}.
\]
This finite space is standard Borel, and \(e\) is jointly measurable. In causal completion zero take \(\mathbf Y=y_0\) at every sample point and define \(Y^{(0)}(a)=e(y_0,a)\) for every \(a\). In causal completion one take \(\mathbf Y=y_1\) at every sample point and define \(Y^{(1)}(a)=e(y_1,a)\) for every \(a\). Thus both completions satisfy the corrected path-typing identity pointwise. In both completions \(\mathbf Y\) is deterministic and therefore conditionally independent of \(A\) given \(X\). Moreover, \(P(A=t_0)=0\), so
\[
Y^{(j)}(A)=e(y_j,A)=0=Y
\quad\text{almost surely},\qquad j\in\{0,1\},
\]
which proves consistency. The observed law is identical, but
\[
\mathbb E[Y^{(0)}(t_0)]=0,
\qquad
\mathbb E[Y^{(1)}(t_0)]=\eta.
\]
For the second completion the counterfactual conditional mean is \(m_P(a,x)=\eta\mathbf 1\{a=t_0\}\), which is discontinuous at \(t_0\). Thus path typing, process-level exchangeability, and consistency alone do not identify the singleton-dose mean; the residual continuity-at-\(t_0\) condition is doing indispensable work.